\newpage\handout
{Part 10: Sequence and Series}
{\textsc{Calculus}}
{\href{https://creativecommons.org/licenses/by-nc-sa/4.0/}{CC BY-NC-SA 4.0 license}}
{Author: \BookAuthor}{Release: \generatedOn}
{Calculus: Part-10}
{\small\textbf{Purpose and Usage:}}\\[0.2cm]
{\footnotesize
These notes represent individual secondary work created by a TA at 
\WorksheetInstitution \space
based on experience with institutional course materials. While derived 
from \WorksheetInstitution \space
content, this compilation includes original explanations, 
examples, and pedagogical approaches developed independently. This work 
is shared as open source educational material and is not officially 
endorsed or authorized by \WorksheetInstitution.}\\[0.5cm]

{\small\textbf{For Educational Use:}}\\[0.2cm]
{\footnotesize
Students and educators are welcome to use and adapt these materials for 
non-commercial educational purposes. Please acknowledge both the original 
\WorksheetInstitution \space
source materials and the individual contributions of the 
author. This work is subject to removal or modification if \WorksheetInstitution 
determines it violates their Honor Code or intellectual property policies.}\\[0.5cm]

{\small\textbf{Content and Attribution:}}\\[0.2cm]
{\footnotesize
This work combines concepts and approaches learned through \WorksheetInstitution 
teaching experience with original pedagogical content. While inspired by 
institutional materials, all explanations, examples, and exercises are 
independently created. This compilation is shared under Creative Commons 
Attribution-NonCommercial-ShareAlike (CC BY-NC-SA) license, with proper 
attribution to both \WorksheetInstitution \space
and the individual author.}\\[0.5cm]

{\small\textbf{Quality Assurance:}}\\[0.2cm]
{\footnotesize
Efforts have been made to ensure clarity and correctness. Users are 
encouraged to verify concepts independently and report any errors or 
concerns. This work is provided as-is for educational benefit and 
may be subject to removal if institutional policies are violated.}

\begin{center}
  \includegraphics[width=0.30\textwidth]{\WorksheetEmblem}\\
  {\small © \the\year\ \BookAuthor. All rights reserved.}\\[0.5cm]
  {\small Prepared by \BookAuthor, TA@GT}\\[1cm]
\end{center}

\newpage


\begin{enumerate}
\item Find the derivatives of the following functions. 
\begin{multicolblock}{In parts (g), (h) and (p) $a$ and $b$ are arbitrary constants.}[3]
  \begin{enumerate}
    \item $y=\left(x^{2}+1\right)e^{3x}$
    \item $y=e^{x^{2}+3x}$
    \item $y=\dfrac{e^{2x}+e^{-2x}}{x^{2}}$
    \item $y=3^{x^{2}+3x}$
    \item $y=x\,5^{3x}$
    \item $y=\dfrac{3^{x}-3^{-x}}{2}$
    \item $y=x e^{a x^{2}+1}$
    \item $y=\sqrt{1+a e^{x}}$
    \item $y=\left(2x+e^{x^{2}}\right)^{4}$
    \item $y=\ln\left(x^{2}+2x\right)$
    \item $y=\log_{2}(3x+4)$
    \item $y=x\ln\left(x^{2}+1\right)$
    \item $y=\log_{3}\left(x^{2}+5\right)$
    \item $y=\dfrac{x\ln x}{x^{2}+1}$
    \item $y=\ln\left(x+5e^{3x}\right)$
    \item $y=a x\ln\left(x^{2}+b^{2}\right)$
    \item $y=(3x)^{5x}$
    \item $y=x^{\ln x}$
    \item $y=(\ln x)^{x}$
    \item $y=(3x+2)^{2x-1}$
  \end{enumerate}
\end{multicolblock}

\item 
\begin{multicolblock}{Solve the equations for $x$:}[3]
  \begin{enumerate}
    \item $2^{x-1}=5$
    \item $3^{2x+3}=7$
    \item $e^{3x+4}=2$
    \item $(3.2)^{x}=64.6$
    \item $\log_{3}(x+4)=1$
    \item $\log_{5}\!\left(x^{2}+9\right)=2$
    \item $\ln(\ln x)=0$
    \item $\ln(x+2)+\ln e^{3}=7$
  \end{enumerate}
\end{multicolblock}

\item This problem deals with functions called the hyperbolic sine and the hyperbolic cosine. These functions occur in the solutions of some differential equations that appear in electromagnetic theory, heat transfer, fluid dynamics, and special relativity. Hyperbolic sine and cosine are defined as follows.
  \[
  \sinh x=\frac{e^{x}-e^{-x}}{2}, \qquad \cosh x=\frac{e^{x}+e^{-x}}{2}.
  \]
  Find derivatives of $\sinh x$ and $\cosh x$ and express your answers in terms of $\sinh x$ and $\cosh x$. Use those formulas to find derivatives of $y=x\sinh x$ and $y=\cosh\left(x^{2}\right)$.


\item For problems 1--6 differentiate the given function.
  \begin{enumerate}
    \item $f(x)=2e^x-8^x$ 
    \item $g(t)=4\log_{3}(t)-\ln(t)$ 
    \item $R(w)=3^w\log(w)$ 
    \item $y=z^5-e^z\ln(z)$ 
    \item $h(y)=\dfrac{y}{1-e^y}$ 
    \item $f(t)=\dfrac{1+5t}{\ln(t)}$ 
  \end{enumerate}

\item Find the tangent line to $f(x)=7^x+4e^x$ at $x=0$. 

\item Find the tangent line to $f(x)=\ln(x)\log_{2}(x)$ at $x=2$. 

\item Determine if $V(t)=\dfrac{t}{e^t}$ is increasing or decreasing at the following points.
  \begin{enumerate}
    \item $t=-4$ 
    \item $t=0$ 
    \item $t=10$ 
  \end{enumerate}

\item Determine if $G(z)=(z-6)\ln(z)$ is increasing or decreasing at the following points.
  \begin{enumerate}
    \item $z=1$ 
    \item $z=5$ 
    \item $z=20$ 
  \end{enumerate}

\end{enumerate}


\newpage






\begin{algorithm}
\caption{Convergence Test of a Series $\sum a_n$}
\label{alg:convergence_test}
\begin{algorithmic}[1]
  \Function{TestConvergence}{$\sum a_n$}
  \If{$\group{\abs{\dfrac{a_n}{b_n}}\nrightarrow \infty}$}
  \Comment{(1) Comparison Test}\State
    \Return \Call{TestConvergence}{$\sum b_n$}
  \ElsIf{$\group{a_n \nrightarrow 0}$} 
  \Comment{(2) Divergence Test}\State
    \Return \inlinequote{Divergent}
  \ElsIf{$\group{\lim_{n \to \infty} \abs{\dfrac{a_{n+1}}{a_n}} < 1}$ 
  \OR $\group{\lim_{n \to \infty} \sqrt[n]{\abs{a_n}} < 1}$} 
  \Comment{(3) Ratio/Root Test}\State
    \Return \inlinequote{Absolutely Convergent}
  \ElsIf{$\group{\abs{a_n} = \dfrac{1}{n^p}}$
  \AND $\group{p > 1}$}
  \Comment{(4) $p$-Series Test}\State
    \Return \inlinequote{Absolutely Convergent}
  \ElsIf{$\group{\int_1^\infty \abs{a_x} \, dx \nrightarrow \infty}$}
  \Comment{(5) Integral Test}\State
    \Return \inlinequote{Absolutely Convergent}
  \ElsIf{$\group{a_n = (-1)^n \abs{a_n}}$}
    \Comment{(6) Alternating Series Test}\State
    \Return \inlinequote{Conditionally Convergent}
  \Else
    \Comment{(7) Failure of Tests}\State
    \Return \inlinequote{Divergent}
  \EndIf
  \EndFunction
\end{algorithmic}
\end{algorithm}









\newpage


\section*{Taylor Series (Maclaurin Series)}
If the function, $f(x)$, and all its derivatives exist at $x=a$ with the 
interval of convergence, then there exists a \textbf{Taylor series} 
(known as \inlinequote{\textbf{polynomial approximation}} or \inlinequote{\textbf{Taylor polynomial}}) 
as follows, denoted $T_{N}(x)$:
$$
T_{N}(x)=\sum_{n=0}^{N} \frac{f^{(n)}(a)}{n!}(x-a)^{n}=f(a)+f^{\prime}(a)(x-a)+\cdots+\frac{f^{(N)}(a)}{N!}(x-a)^{N} \quad \text { where } \quad|x-a|<R
$$
where
\begin{enumerate}
\item $(a-R, a+R)$ is the \textbf{interval of convergence}.
\item $R>0$ is the \textbf{radius of convergence}.
\item $N$ is the \textbf{degree of Taylor series} (or up-to \textbf{exponent $N$}).
\end{enumerate}
\textbf{Note}: Maclaurin series is specifically a Taylor series at $x=0$. It is very important since it lets us easily substitute, integrate, and differentiate to find representations of other functions as follows.

\textbf{Example}. The following is a Maclaurin series of rational function,
$$
\frac{1}{1-x}=\sum_{n=0}^{\infty} x^{n} \quad \text { where } \quad|x|<1
$$
Using this, any rational functions can be derived into series as well as its interval of convergence. The following are examples of the application.
$$
\begin{aligned}
& \frac{2}{3-x}=\frac{2}{3} \cdot \frac{1}{1-\left(\frac{x}{3}\right)}=\frac{2}{3} \cdot \sum_{n=0}^{\infty}\left(\frac{x}{3}\right)^{n} \quad \text { where } \quad\left|\frac{x}{3}\right|<1 \\
& \frac{1}{x}=\frac{1}{1+(x-1)}=\sum_{n=0}^{\infty}(x-1)^{n} \quad \text { where }|x-1|<1
\end{aligned}
$$
Also, the following is a Maclaurin series of trigonometric function,
$$
\cos x=\sum_{n=0}^{\infty}(-1)^{n} \frac{x^{2 n}}{(2 n)!}
$$
By one of the trigonometric properties,
$$
\sin \left(x+\frac{\pi}{2}\right)=\cos x
$$
Thus,
$$
\sin x=\cos \left(x-\frac{\pi}{2}\right)=\sum_{n=0}^{\infty}(-1)^{n} \frac{\left(x-\frac{\pi}{2}\right)^{2 n}}{(2 n)!}
$$
This shows how to get the Taylor series of sine at $x=\frac{\pi}{2}$. With this, time will be saved a lot. (No need to get every single $f^{(n)}(a)$, all the derivatives of $f$ )


\newpage
\section*{Maclaurin Series (Taylor Series at $x=0$)}
Geometric Series (Rational and Logarithmic function)
$$
\begin{aligned}
\frac{1}{1-x} & =\sum_{n=0}^{\infty} x^{n}=1+x+x^{2}+x^{3}+\cdots & & \text { where }|x|<1 \\
\ln (1+x) & =\sum_{n=0}^{\infty}(-1)^{n} \frac{x^{n+1}}{n+1}=x-\frac{x^{2}}{2}+\frac{x^{3}}{3}-\frac{x^{4}}{4}+\cdots & & \text { where }|x|<1
\end{aligned}
$$
Binomial Series (Power function, general\footnote{(Undergraduate Level) The Binomial Series is considered 
\inlinequote{general} because it applies to any real or complex 
exponent, not just to integers. It provides a flexible way to get 
any rational and logarithmic series.})
$$
(1+x)^{p}=\sum_{n=0}^{\infty} \frac{p(p-1) \cdots(p-n+1)}{n!} x^{n} \quad \text { where }|x|<1
$$
Transcendental Series (Exponential and Trigonometric function)
$$
\begin{aligned}
\cos x & =\sum_{n=0}^{\infty}(-1)^{n} \frac{x^{2 n}}{(2 n)!}=1-\frac{x^{2}}{2!}+\frac{x^{4}}{4!}-\frac{x^{6}}{6!}+\cdots & & \text { where } x \in \mathbb{R} \\
\sin x & =\sum_{n=0}^{\infty}(-1)^{n} \frac{x^{2 n+1}}{(2 n+1)!}=x-\frac{x^{3}}{3!}+\frac{x^{5}}{5!}-\frac{x^{7}}{7!}+\cdots & & \text { where } x \in \mathbb{R} \\
e^{x} & =\sum_{n=0}^{\infty} \frac{x^{n}}{n!}=1+x+\frac{x^{2}}{2!}+\frac{x^{3}}{3!}+\frac{x^{4}}{4!}+\cdots & & \text { where } x \in \mathbb{R}
\end{aligned}
$$
Euler's formula (Complex function, general\footnote{(Undergraduate Level) Euler's formula bridges the 
gap between exponential functions, trigonometric functions, and 
complex numbers, revealing their interrelationships and demonstrating 
the unity of mathematical concepts across different domains. It provides 
a powerful tool for understanding and analyzing complex phenomena in 
mathematics, physics, engineering, and other disciplines.})
$$
e^{i x}=\cos (x)+i \cdot \sin (x) \quad \text { where } \quad x \in \mathbb{R}
$$




\newpage

\section*{Examples with Solutions}
Find the Taylor polynomial of degree 3 for

$$
f(x)=\frac{1}{x+2}
$$
about the point $x=3$.\\
\begin{subanswer}
\begin{enumerate}
\item (\textbf{Taylor series}). By definition,
$$
T_{3}(x)=\sum_{n=0}^{3} \frac{f^{(n)}(3)}{n!}(x-3)^{n}=\frac{f(3)}{0!}(x-3)^{0}+\frac{f^{\prime}(3)}{1!}(x-3)^{1}+\frac{f^{\prime \prime}(3)}{2!}(x-3)^{2}+\frac{f^{\prime \prime \prime}(3)}{3!}(x-3)^{3}
$$
\begin{enumerate}
  \item find all the derivatives.
  $$
  \begin{array}{rlrl}
  f(3) & =\dfrac {1}{5} & & \\
  f^{\prime}(x) & =\dfrac{-1}{(x+2)^{2}} \text {, that is, } \quad f^{\prime}(3)=\dfrac {-1}{25} \\
  f^{\prime \prime}(x) & =\dfrac {2}{(x+2)^{3}} \text {, that is, } \quad f^{\prime \prime}(3)=\dfrac {2}{125} \\
  f^{\prime \prime \prime}(x) & =\dfrac {-6}{(x+2)^{4}} \text {, that is, } \quad f^{\prime \prime \prime}(x)=\dfrac {-6}{625}
  \end{array}
  $$
  \item find the coefficients.
  $$
  \begin{aligned}
  \frac{f^{\prime}(3)}{1!} & =\frac{-\frac{1}{25}}{1}=\frac{-1}{25} \\
  \frac{f^{\prime \prime}(3)}{2!} & =\frac{\frac{2}{125}}{2}=\frac{1}{125} \\
  \frac{f^{\prime \prime \prime}(3)}{3!} & =\frac{6 / 625}{6}=\frac{1}{125}
  \end{aligned}
  $$
  \item apply them to the Taylor polynomial.
  $$
  T_{3}(x)=\frac{1}{5}-\frac{(x-3)}{25}+\frac{(x-3)^{2}}{125}-\frac{(x-3)^{3}}{625}
  $$
\end{enumerate}
Fortunately, it gives the right answer. However, this approach was painful in terms of calculation labor. Instead, the following answer can be a good alternative.
\item (\textbf{Maclaurin series}). We know the Maclaurin series already. So, with some algebra,
$$
f(x)=\frac{1}{x+2}=\frac{1}{5} \cdot \frac{1}{1-\left(-\frac{x-3}{5}\right)}=\frac{1}{5} \cdot \sum_{n=0}^{\infty}(-1)^{n} \cdot\left(\frac{x-3}{5}\right)^{n} \quad \text { where } \quad\left|\frac{x-3}{5}\right|<1
$$
Therefore,
$$
T_{3}(x)=\frac{1}{5} \cdot \sum_{n=0}^{\infty}(-1)^{n} \cdot\left(\frac{x-3}{5}\right)^{n}=\frac{1}{5}-\frac{1}{25}(x-3)+\frac{1}{125}(x-3)^{2}-\frac{1}{625}(x-3)^{3}
$$
\end{enumerate}
\end{subanswer}



\newpage 


\section*{Extra Examples with Short Answers}

Find the Maclaurin series for $f(x)=\frac{1}{1+x^{2}}$.
\begin{subanswer}
Begin with the known series.
$$
f(x)=\frac{1}{1-x}=\sum_{n=0}^{\infty} x^{n}
$$

Substitute $-x^{2}$ for $x$.
$$
f(x)=\frac{1}{1+x^{2}}=\frac{1}{1-\left(-x^{2}\right)}=\sum_{n=0}^{\infty}\left(-x^{2}\right)^{n}=\sum_{n=0}^{\infty}(-1)^{n} x^{2 n}=1-x^{2}+x^{4}-x^{6}+\cdots
$$
\end{subanswer}

A function $f(x)$ is approximated by the third order Taylor series $1+2(x-1)-(x-1)^{2}+(x-1)^{3}$ centered at $x=1$. Find $f^{\prime}(1)$ and $f^{\prime \prime}(1)$.
\begin{subanswer}
The question can also be asking for the derivative's values reversely. In this case, there is no such painful calculation labor. Thus, to use the definition of Taylor series should be good enough as follows.

Compare $\sum_{n=0}^{\infty} \frac{f^{(n)}(a)}{n!}(x-a)^{n}$ to the given polynomial:
$$
\begin{gathered}
\frac{f(1)}{0!}=1 \\
\frac{f^{\prime}(1)}{1!}=2 \\
\frac{f^{\prime \prime}(1)}{2!}=-1 \\
\vdots
\end{gathered}
$$

Solve for derivatives as needed.
$$
\begin{gathered}
f^{\prime}(1)=2 \cdot 1!=2 \\
f^{\prime \prime}(1)=-1 \cdot 2!=-2
\end{gathered}
$$
\end{subanswer}

Differentiate the Maclaurin series for $f(x)=\ln (x+1)$ to find the Taylor series expansion for $f(x)=\frac{1}{x+1}$.
\begin{subanswer}
The Taylor series is also handy for dealing with differentiation and integration.\\
Use the existing Maclaurin series formula.
$$
f(x)=\ln (x+1)=x-\frac{1}{2} x^{2}+\frac{1}{3} x^{3}-\frac{1}{4} x^{4}+\cdots
$$
Find derivatives over summation ($\sum$).
$$
f^{\prime}(x)=\frac{1}{x+1}=1-x+x^{2}-x^{3}+\cdots=\sum_{n=0}^{\infty}(-1)^{n} x^{n}
$$
\end{subanswer}














\newpage
\section*{Practices}
\begin{enumerate}
\item (Use definition of Taylor series) Find the Maclaurin polynomial of degree 4 that approximates $f(x)=\ln (1+x)$.
\item (Try both methods) Find the Taylor series for the function $f(x)=e^{-x}$ about the point $x=\ln 2$.
\item Find the Maclaurin series for the function $f(x)=x e^{x}$.
\item Find the first four non-zero terms of the Maclaurin series for $f(x)=x \cdot \cos (2 x)$.
\item Approximate the function with a fourth degree Taylor polynomial centered at $x=1$.
$$
f(x)=e^{x^{2}}
$$
\item Find the Maclaurin series for the function and determine its interval of convergence.
$$
f(x)=\frac{1}{1+x^{2}}
$$

\item Use a Maclaurin series to approximate the integral
$$
\int_{0}^{1} \sin \left(x^{2}\right) d x
$$
\end{enumerate}
\begin{multicolblock}{Determine whether the 
following series converge or diverge. 
In the case of convergence, give the sum of the series.}[3]
\begin{enumerate}
  \item \[\sum_{n=1}^{\infty} \frac{1}{2 n}\]
  \item \[\sum_{n=1}^{\infty} \frac{n+1}{2 n-3}\]
  \item \[\sum_{n=2}^{\infty} \frac{n^{2}}{n^{2}-1}\]
  \item \[\sum_{n=1}^{\infty} \frac{n(n+2)}{(n+3)^{2}}\]
  \item \[\sum_{n=1}^{\infty} \frac{1+2^{n}}{3^{n}}\]
  \item \[\sum_{n=1}^{\infty} \frac{1+3^{n}}{2^{n}}\]
  \item \[\sum_{n=1}^{\infty} \sqrt[n]{2}\]
  \item \[\sum_{n=1}^{\infty} 0.8^{n-1}-0.3^{n}\]
  \item \[\sum_{n=1}^{\infty} \ln \left(\frac{n^{2}+1}{2 n^{2}+1}\right)\]
  \item \[\sum_{n=1}^{\infty} \cos ^{n}(1)\]
  \item \[\sum_{n=1}^{\infty} \tan ^{-1}(n)\]
  \item \[\sum_{n=1}^{\infty}\left(\frac{3}{5^{n}}+\frac{2}{n}\right)\]
  \item \[\sum_{n=1}^{\infty}\left(\frac{1}{e^{n}}+\frac{1}{n(n+1)}\right)\]
  \item \[\sum_{n=1}^{\infty} \frac{e^{n}}{n^{2}}\]
  \item \[\sum_{n=2}^{\infty} \frac{2}{n^{2}-1}\]
  \item \[\sum_{n=1}^{\infty} \frac{2}{n^{2}+4 n+3}\]
  \item \[\sum_{n=1}^{\infty} \frac{3}{n(n+3)}\]
  \item \[\sum_{n=1}^{\infty} \ln \left(\frac{n}{n+1}\right)\]
  \item \[\sum_{n=1}^{\infty}\left(e^{1 / n}-e^{1 /(n+1)}\right)\]
  \item \[\sum_{n=1}^{\infty} 6(0.9)^{n-1}\]
\end{enumerate}
\end{multicolblock}
\newpage
\begin{multicolblock}{Write the following numbers as a ratio of two integers.}[3]
\begin{enumerate}
  \item $22$
  \item $0.222222 \ldots$
  \item $243.417417 \ldots$
  \item $261.53424242 \ldots$
  \item $230.737373 \ldots$
  \item $256.2545454 \ldots$
  \item $7.1234512345 \ldots$
\end{enumerate}
\end{multicolblock}